% ============================================================
\section{Thiết kế và hiện thực}
\label{sec:arch}
% ============================================================

\subsection{Kiến trúc}

Hệ thống là một pipeline sáu bước, mỗi bước biến đổi câu hỏi sang một biểu diễn mới: dãy
token, cây cú pháp, frame ngữ nghĩa, frame đã giải tham chiếu, dữ kiện và câu trả lời, như ở
\cref{fig:pipeline}. Ý tưởng chính là tách việc \emph{hiểu câu hỏi} khỏi việc \emph{tra dữ
kiện}: tầng ngữ nghĩa chỉ dựa vào cây cú pháp và từ điển thực thể, không đọc KB. Nhờ vậy hệ
thống phân biệt được câu không hiểu với câu hiểu nhưng KB không có dữ kiện, và câu trả lời
chỉ chứa những gì truy vấn tìm thấy.

\begin{figure}[tbp]
  \centering
  \begin{tikzpicture}[
      font=\small,
      stage/.style={draw, rounded corners=2pt, fill=blue!6, minimum width=6cm,
        minimum height=8.5mm, align=center},
      res/.style={draw, rounded corners=2pt, fill=orange!10, minimum width=3.4cm,
        minimum height=8.5mm, align=center, font=\footnotesize},
      io/.style={draw, rounded corners=6pt, fill=green!8, minimum width=6cm,
        minimum height=7mm, align=center},
      arr/.style={-{Stealth[length=2mm]}, thick},
      lab/.style={font=\scriptsize\itshape, text=black!65, right=1mm}
    ]
    \node[io] (in) {Câu hỏi tiếng Việt};
    \node[stage, below=6mm of in] (tok) {Tách token \\ \scriptsize chuẩn hóa Unicode, chữ thường, tách dấu câu};
    \node[stage, below=6mm of tok] (parse) {Parser Earley \\ \scriptsize nhận hoặc từ chối câu, dựng cây};
    \node[stage, below=6mm of parse] (sem) {Diễn giải ngữ nghĩa \\ \scriptsize nhánh câu hỏi $\rightarrow$ intent, slot};
    \node[stage, below=6mm of sem] (ctx) {Giải tham chiếu \\ \scriptsize thay ``nó'', ``chương này'' bằng thực thể đã nhớ};
    \node[stage, below=6mm of ctx] (query) {Truy vấn KB \\ \scriptsize tra cứu theo intent và slot};
    \node[stage, below=6mm of query] (ans) {Sinh câu trả lời \\ \scriptsize template theo loại dữ kiện};
    \node[io, below=6mm of ans] (outn) {Câu trả lời};

    \node[res, left=9mm of parse] (cfg) {Văn phạm CFG \\ 26 nhánh câu hỏi};
    \node[res, left=9mm of sem] (lex) {Từ điển thực thể};
    \node[res, left=9mm of query] (kb) {Knowledge base};

    \foreach \a/\b/\t in {in/tok/{}, tok/parse/{token}, parse/sem/{cây cú pháp},
        sem/ctx/{frame}, ctx/query/{frame đã giải}, query/ans/{dữ kiện}, ans/outn/{}}
      {\draw[arr] (\a) -- (\b) node[midway, lab] {\t};}
    \draw[arr, dashed] (cfg) -- (parse);
    \draw[arr, dashed] (lex) -- (sem);
    \draw[arr, dashed] (kb) -- (query);
  \end{tikzpicture}
  \caption{Pipeline xử lý một câu hỏi}
  \label{fig:pipeline}
\end{figure}

\subsection{Văn phạm (Phần I)}

Trước khi phân tích, câu được chuẩn hóa Unicode, chuyển chữ thường và tách dấu câu, nhưng
giữ nguyên các thuật ngữ như ``context-free'' hay ``LO2.5''.

Văn phạm gồm 100 ký hiệu không kết thúc, 463 luật sinh và 330 terminal, viết theo kiểu
\emph{semantic grammar}: ký hiệu không kết thúc mang nghĩa của miền như \nt{WEEK\_ENTITY}
hay \nt{Q\_SCHEDULE\_WEEK} thay vì nhãn từ loại như NP, VP. Mỗi câu hợp lệ đi qua đúng một
trong 26 nhánh câu hỏi, và nhánh đó quyết định intent; \cref{tab:branches} cho biết nhánh
nào ứng với intent nào. Thực thể nằm sẵn trong các node có tên, nên bước diễn giải chỉ cần
đọc cây. Ví dụ một phần văn phạm cho câu hỏi lịch học theo tuần:
\begin{center}
  \begin{tabular}{r@{\ $\rightarrow$\ }l}
    \nt{S} & \nt{PREFIX} \nt{QUESTION} \nt{SUFFIX} \\
    \nt{QUESTION} & \nt{Q\_COURSE\_CREDITS} $\mid$ \nt{Q\_SCHEDULE\_WEEK} $\mid$ \dots \\
    \nt{Q\_SCHEDULE\_WEEK} & \nt{WEEK\_ENTITY} \nt{WEEK\_CONTENT\_QUESTION} $\mid$ \dots \\
    \nt{WEEK\_ENTITY} & \textit{tuần} \nt{WEEK\_NUMBER} \\
    \nt{WEEK\_NUMBER} & \textit{1} $\mid$ \textit{2} $\mid$ \dots $\mid$ \textit{16}
      $\mid$ \textit{20} $\mid$ \textit{99} \\
  \end{tabular}
\end{center}
\nt{WEEK\_NUMBER} cố ý nhận cả tuần không có trong lịch như tuần 99: câu đúng cú pháp thì
được hiểu, còn tuần đó có tồn tại hay không là việc của KB. Vì vậy ``Tuần 99 học gì?'' được
trả lời là không tìm thấy thay vì bị coi là câu ngoài văn phạm.

\begin{table}[htbp]
  \centering
  \small
  \begin{tabularx}{\linewidth}{|P{3.3cm}|L|}
    \hline
    \textbf{Intent} & \textbf{Nhánh câu hỏi trong văn phạm} \\ \hline
    \code{GET_COURSE_INFO} & \code{Q_COURSE_CREDITS}, \code{Q_COURSE_PARTS},
      \code{Q_COURSE_DESCRIPTION}, \code{Q_COURSE_EXAM_DURATION},
      \code{Q_COURSE_EXAM_TIME}, \code{Q_COURSE_ASSESSMENT} \\ \hline
    \code{GET_INSTRUCTOR} & \code{Q_INSTRUCTOR} \\ \hline
    \code{GET_SCHEDULE} & \code{Q_SCHEDULE_WEEK}, \code{Q_SCHEDULE_CHAPTER},
      \code{Q_SCHEDULE_TOPICS} \\ \hline
    \code{GET_TOPIC} & \code{Q_TOPIC_BY_NAME}, \code{Q_TOPIC_BY_CHAPTER},
      \code{Q_TOPIC_RELATION} \\ \hline
    \code{GET_LO} & \code{Q_LO} \\ \hline
    \code{GET_ASSIGNMENT} & \code{Q_ASSIGNMENT_OVERVIEW}, \code{Q_ASSIGNMENT_PART},
      \code{Q_ASSIGNMENT_FEATURE}, \code{Q_ASSIGNMENT_OUTPUT} \\ \hline
    \code{GET_DEADLINE} & \code{Q_DEADLINE} \\ \hline
    \code{GET_RESOURCE} & \code{Q_RESOURCE} \\ \hline
    \code{GET_RULE} & \code{Q_RULE} \\ \hline
    \code{CONTEXT_DEPENDENT} & \code{Q_CONTEXT_TOPIC_RELATION},
      \code{Q_CONTEXT_NEXT_CHAPTER}, \code{Q_CONTEXT_ASSIGNMENT_OUTPUT},
      \code{Q_CONTEXT_RESOURCE}, \code{Q_CONTEXT_LO_CHAPTER} \\ \hline
    \code{UNKNOWN} & Không có cây \\ \hline
  \end{tabularx}
  \caption{26 nhánh câu hỏi và intent tương ứng}
  \label{tab:branches}
\end{table}

Ở giai đoạn 1, mỗi thành viên viết một nhóm nhánh; bản gộp giữ nguyên 26 nhánh và loại các
lựa chọn gây mơ hồ. Các cách gọi đồng nghĩa được gom vào cùng một ký hiệu, ví dụ
``parser'', ``parse tree'', ``phân tích cú pháp'' đều thuộc \nt{FEAT\_PARSING} và cùng chỉ
Phần I của bài tập lớn.

\subsection{Parser Earley và sinh câu (Phần I)}
\label{sec:grammarcheck}

Nhóm dùng thuật toán Earley \cite{earley1970,jurafsky} vì nó chạy được trên mọi CFG, kể cả
luật rỗng và đệ quy trái, mà không phải đổi văn phạm sang dạng chuẩn Chomsky như CKY, nên
cây giữ nguyên các node ngữ nghĩa. Với $n$ token, parser dựng các tập trạng thái
$S_0, \dots, S_n$ bằng ba thao tác:
\begin{itemize}
  \item \textbf{Predictor}: gặp ký hiệu không kết thúc $B$ sau dấu chấm, thêm mọi luật
        $B \rightarrow \bullet\,\gamma$.
  \item \textbf{Scanner}: terminal sau dấu chấm khớp token kế tiếp, đưa trạng thái sang
        $S_{i+1}$.
  \item \textbf{Completer}: $B \rightarrow \gamma\,\bullet$ hoàn tất, dịch dấu chấm cho các
        trạng thái đang chờ $B$ và gắn cây con.
\end{itemize}
Câu được chấp nhận khi $S_n$ có \nt{S} hoàn tất từ vị trí 0; câu bị từ chối nhận intent
\code{UNKNOWN}. Parser tìm tối đa hai cây để phát hiện câu mơ hồ. Độ phức tạp xấu nhất là
$O(n^3)$.

\begin{figure}[tbp]
  \centering
  \begin{adjustbox}{max width=\linewidth}
  \begin{forest}
    for tree={font=\footnotesize\sffamily, s sep=4mm, l sep=4mm, inner sep=1.5pt,
      edge={black!60}},
    where n children=0{font=\small\itshape, text=blue!50!black}{}
    [S
      [PREFIX [$\varepsilon$]]
      [QUESTION
        [Q\_SCHEDULE\_WEEK
          [WEEK\_ENTITY [tuần] [WEEK\_NUMBER [3]]]
          [WEEK\_CONTENT\_QUESTION [học] [gì]]]]
      [SUFFIX [END\_MARK [?]]]]
  \end{forest}
  \end{adjustbox}
  \caption{Cây cú pháp của câu ``Tuần 3 học gì?''}
  \label{fig:tree}
\end{figure}

Câu được sinh bằng dẫn xuất trái nhất theo chiều rộng, xoay vòng giữa 26 nhánh để nhánh nào
cũng có câu. Dạng câu chắc chắn dài quá 32 token bị cắt sớm, nên quá trình luôn dừng kể cả
với luật đệ quy và luật rỗng. Mọi câu sinh ra đều được parse lại để chắc chắn thuộc văn phạm.

Sau mỗi lần sửa văn phạm, nhóm kiểm tra luật trùng và ký hiệu không dùng tới, rồi chạy lại
toàn bộ câu kiểm thử và 1.000 câu sinh. Ở lần chạy gần nhất, 1.000 câu sinh phủ đủ 26 nhánh
và đều parse lại được; không câu nào trong 1.047 câu được kiểm tra có hai cây.

\subsection{Ngữ nghĩa, truy vấn và trả lời (Phần II)}

Ngữ nghĩa được dựng theo nguyên tắc tính hợp thành: nghĩa của câu ghép từ nghĩa các thành
phần có tên trên cây. Nhánh câu hỏi cho intent, các node con như \nt{WEEK\_ENTITY} cho slot,
và từ điển thực thể chuẩn hóa cụm từ về ID của KB, ví dụ ``tuần 3'' $\rightarrow$ WEEK\_03,
``NLP'' $\rightarrow$ CO3085. Kết quả là một frame dạng vị từ
\textit{intent}(\textit{slot} = \textit{giá trị}), xem \cref{tab:frames}.

\begin{table}[htbp]
  \centering
  \small
  \begin{tabularx}{\linewidth}{|P{4cm}|L|}
    \hline
    \textbf{Câu hỏi} & \textbf{Frame ngữ nghĩa} \\ \hline
    Môn NLP có bao nhiêu tín chỉ? & \code{GET_COURSE_INFO(course=CO3085, field=CREDITS)} \\ \hline
    Tuần 3 học gì? & \code{GET_SCHEDULE(week=WEEK_03)} \\ \hline
    BTL có phần parse tree không? & \code{GET_ASSIGNMENT(section=PART_I, feature=parse tree)} \\ \hline
    Có được dùng thư viện có sẵn không? & \code{GET_RULE(rule=EXTERNAL_LIBRARIES)} \\ \hline
    Hôm nay trời có mưa không? & \code{UNKNOWN()} \\ \hline
  \end{tabularx}
  \caption{Ví dụ frame ngữ nghĩa}
  \label{tab:frames}
\end{table}

Mỗi intent ứng với một cách tra KB. Tuần, chương, LO và quy định được tra chính xác theo ID;
chủ đề và tài liệu được tra theo từ khóa, ưu tiên khớp nguyên cụm. Câu trả lời được điền vào
template theo loại dữ kiện. Ví dụ, ``Môn NLP có bao nhiêu tín chỉ?'' đi qua nhánh
\nt{Q\_COURSE\_CREDITS}, cho frame hỏi trường tín chỉ của môn CO3085, và nhận câu trả lời
``Môn Xử lý Ngôn ngữ Tự nhiên (CO3085) có 3 tín chỉ.''

Khi không trả lời được, hệ thống tách ba trường hợp ở \cref{tab:fallback}, để người dùng biết
nên hỏi lại, bổ sung ngữ cảnh, hay dữ kiện thực sự không có. Câu cuối là nguyên văn thông báo
đề yêu cầu \cite{spec}.

\begin{table}[htbp]
  \centering
  \small
  \begin{tabularx}{\linewidth}{|P{3.2cm}|P{3.6cm}|L|}
    \hline
    \textbf{Trường hợp} & \textbf{Dấu hiệu} & \textbf{Câu trả lời} \\ \hline
    Ngoài văn phạm & Không có cây cú pháp &
      Xin lỗi, câu hỏi không thuộc văn phạm hoặc ngoài phạm vi Course Assistant. \\ \hline
    Thiếu ngữ cảnh & Câu tham chiếu nhưng chưa có thực thể để thay &
      Xin lỗi, tôi chưa có đủ ngữ cảnh để hiểu tham chiếu trong câu hỏi này. \\ \hline
    KB thiếu dữ kiện & Hiểu câu nhưng truy vấn rỗng &
      Xin lỗi, tôi không tìm thấy thông tin này trong dữ liệu của môn học. \\ \hline
  \end{tabularx}
  \caption{Ba loại fallback}
  \label{tab:fallback}
\end{table}

\subsection{Ngữ cảnh hội thoại (Phần III)}

Phần III xử lý tham chiếu cục bộ: câu hỏi được dùng ``nó'', ``chương này'', ``phần đó'',
``chương tiếp theo'' để chỉ thực thể vừa nhắc. Văn phạm có năm nhánh cho năm dạng tham
chiếu; các câu này nhận intent \code{CONTEXT_DEPENDENT} và được giải trước khi truy vấn theo
ba nguyên tắc:
\begin{itemize}
  \item nhớ thực thể gần nhất riêng cho từng loại: chương, phần bài tập lớn, LO, để câu hỏi
        về bài tập lớn không xóa chương đang nhớ;
  \item thay tham chiếu bằng thực thể gần nhất đúng loại, nếu chưa có thì báo thiếu ngữ
        cảnh;
  \item chỉ ghi nhớ thực thể khi truy vấn tìm thấy dữ kiện, nên câu hỏi về chương 13 không
        thay chương đang nhớ.
\end{itemize}
Ví dụ, sau ``Tuần 3 học gì?'' hệ thống nhớ Chương 3; câu ``Nó có liên quan đến parsing
không?'' được giải ``nó'' thành Chương 3, tra chủ đề của chương và trả lời ``Có''. Cả 6 hội
thoại mẫu của đề đều giải đúng.
